\uuid{kkjR}
\exo7id{4180}
\titre{exo7 4180}
\auteur{quercia}
\organisation{exo7}
\datecreate{2010-03-11}
\isIndication{false}
\isCorrection{false}
\chapitre{Fonction de plusieurs variables}
\sousChapitre{Dérivée partielle}
\cpge{mpsi-mp2i.fonctions-deux-variables.b, pcsi.fonctions-deux-variables.b, ptsi.fonctions-deux-variables.b, mpsi-mp2i.fonctions-deux-variables.c, pcsi.fonctions-deux-variables.c, ptsi.fonctions-deux-variables.c, bcpst1.fonctions-deux-variables.c, mp-mpi.calcul-differentiel-optimisation.b, pc.fonctions-plusieurs-variables.c, psi.fonctions-plusieurs-variables.c}
\module{Analyse}
\niveau{L2}
\difficulte{}

\contenu{
\question{
Soit $f:\R \to \R$ de classe $\mathcal{C}^1$ et $k$-lipschitzienne avec $k<1$ et
$\varphi : {\R^2} \to {\R^2}, {(x,y)} \mapsto {(x+f(y),y-f(x)).}$

Montrer que $\varphi$ est un $\mathcal{C}^1$-difféomorphisme de~$\R^2$ sur~$\R^2$.
}
}
